\subsection{Matricies and Matrix Operations}
\hrulefill

\paragraph*{Definition} A \textbf{matrix} is a rectangular array of numbers arranged in rows and columns. An $m \times n$ matrix has $m$ rows and $n$ columns. For example, the following is a $2 \times 3$ matrix:
\[A = \begin{bmatrix}
    1 & 2 & 3 \\
    4 & 5 & 6
\end{bmatrix}\]

\paragraph*{Matrix Access} The entry in the $i$-th row and $j$-th column of a matrix $A$ is denoted by $a_{ij}$. For example, in the matrix $A$ above, $a_{12} = 2$ and $a_{21} = 4$.

\[A_{m \times n} = \begin{bmatrix}
    a_{11} & a_{12} & \ldots & a_{1n} \\
    a_{21} & a_{22} & \ldots & a_{2n} \\
    \vdots & \vdots & \ddots & \vdots \\
    a_{m1} & a_{m2} & \ldots & a_{mn}
\end{bmatrix}\]

\subsubsection*{Matrix Operations}

\paragraph*{Matrix Equality} Two matrices $A$ and $B$ are equal if they have the same dimensions and their corresponding entries are equal. That is, $A = B$ if and only if $a_{ij} = b_{ij}$ for all $i$ and $j$.

\paragraph*{Matrix Addition} Two matrices can be added if they have the same dimensions, otherwise it is undefined. If $A$ and $B$ are both $m \times n$ matrices, then their sum is given by
\[A + B = \begin{bmatrix}
    a_{11} + b_{11} & a_{12} + b_{12} & \ldots & a_{1n} + b_{1n} \\
    a_{21} + b_{21} & a_{22} + b_{22} & \ldots & a_{2n} + b_{2n} \\
    \vdots & \vdots & \ddots & \vdots \\
    a_{m1} + b_{m1} & a_{m2} + b_{m2} & \ldots & a_{mn} + b_{mn}
\end{bmatrix}\]

\paragraph*{Scalar Multiplication} If $A$ is an $m \times n$ matrix and $c$ is a scalar, then the scalar multiplication of $c$ and $A$ is given by
\[cA = \begin{bmatrix}
    ca_{11} & ca_{12} & \ldots & ca_{1n} \\
    ca_{21} & ca_{22} & \ldots & ca_{2n} \\
    \vdots & \vdots & \ddots & \vdots \\
    ca_{m1} & ca_{m2} & \ldots & ca_{mn}
\end{bmatrix}\]

\paragraph*{Properties of Matrix Addition and Scalar Multiplication}
Let $A$, $B$, and $C$ be matrices of the same dimensions, and let $c$ and $d$ be scalars. Then the following properties hold:
\begin{itemize}
    \item Commutative Property: $A + B = B + A$
    \item Associative Property: $(A + B) + C = A + (B + C)$
    \item Distributive Property: $c(A + B) = cA + cB$
    \item Scalar Distributive Property: $(c + d)A = cA + dA$
    \item Associative Property of Scalar Multiplication: $c(dA) = (cd)A$
    \item Identity Property: $A + 0 = A$, where $0$ is the zero matrix of the same dimensions as $A$
\end{itemize}

\paragraph*{Matrix Transpose} The transpose of a matrix $A$, denoted by $A^T$, is obtained by interchanging its rows and columns. If $A$ is an $m \times n$ matrix, then $A^T$ is an $n \times m$ matrix given by
\[A^T = \begin{bmatrix}
    a_{11} & a_{21} & \ldots & a_{m1} \\
    a_{12} & a_{22} & \ldots & a_{m2} \\
    \vdots & \vdots & \ddots & \vdots \\
    a_{1n} & a_{2n} & \ldots & a_{mn}
\end{bmatrix}\]

\paragraph*{Transposition Identities}
Let $A, B$ be $m\times n$ matricies and $k \in \mathbb{R}$
\begin{itemize}
    \item $A^T$ is $n \times m$
    \item $(A^T)^T = A$
    \item $ (kA)^T = k(A^T)$
    \item $(A + B)^T = A^T + B^T$
\end{itemize}

\paragraph*{Example}
Solve for A
\begin{align*}
    (2A^T + 3 \begin{bmatrix}
        1 & 2 \\
        -1 & 1
    \end{bmatrix})^T
\end{align*}

\paragraph*{Symmetric Matrix} A matrix $A$ is called symmetric if $A = A^T$. If A, B are $n \times n$ symmetric matrices, then $A + B$ is also symmetric. If c is a scalar, then cA is also symmetric.
If $A$ and $B$ are both symmetric, then $kA + kB$ is also symmetric for $k, p \in \mathbb{R}$.